\Russian

%%%%

%%%   Самарский государственный технический университет

%%%   Кафедра прикладной математики и информатики

%%%

%%%   Преамбула для статьи в журнал "Вестник Самарского государственного

%%%   технического университета. Серия: «Физико-математические науки»"

%%%   Ориентирована под MikTEx 2.4 for Win32

%%%

%%%   Хотя сам журнал собирается под GNU/Linux на дистрибутиве TeXLive



\documentclass[11pt,a4paper,twoside]{article}



\usepackage[cp1251]{inputenc}

\usepackage[T2A]{fontenc}

\usepackage[english,russian]{babel}

\usepackage{samgtubib}

\usepackage[dvips]{graphicx}

\usepackage[multiple]{footmisc}



\usepackage{samgtu}

\renewcommand{\VestnikYear}{2009} % Год издания Вестника

\renewcommand{\VestnikNumber}{2\,(19)} % Номер Вестника

\renewcommand{\VestnikMonth}{Ноябрь} % Месяц выхода Вестника

\renewcommand{\VestnikData}{01.11.09} % Дата подписания Вестника в печать

\renewcommand{\VestnikUPL}{26,64} % Усл. печ. л.

\renewcommand{\VestnikUIL}{25,73} % Уч.-изд. л.

\renewcommand{\VestnikRegNumber}{479/09} % Регистрационный номер

\renewcommand{\VestnikOrder}{994} % Номер заказа






%
%%\samgtupubl % для генерации pdf, отдаваемого в типографию СамГТУ
%\pdfcrop % обрезка pdf для размещения на math-net.ru
%\draft % На корректуре
\nofilllastpage % Последняя страница не заполнена не полностью
%\briefcomm % Статья является кратким сообщением


\vsgtu{1265} % Номер статьи по базе Math-Net.ru

\FirstPage{19}
\LastPage{24}


\renewcommand{\baselinestretch}{0.95}

%\begin{document}


% Номер УДК
\udk{517.442+517.581}%Обработка текста. Подготовка текстов : Языки программирования. Компьютерные языки
\msc{44A10, 44A20; 33C15, 33C20}% Коды MSC см. http://www.ams.org/msc/

\rutitle{Обобщ\"eнное интегральное преобразование Лапласа и~его применение к~решению некоторых интегральных уравнений}% Полный заголовок
{Обобщённое интегральное преобразование Лапласа \dots}% Заголовок, который идёт в верхний колонтитул

\ruauthor{С.~М.~Заикина\inst{1,2}}{С.~М.~З\,а\,и\,к\,и\,н\,а}

\ruaffil{\rusamgtu. \and \ruvolgogradsu.}

\email{1}{\mem{svetzai@inbox.ru}} %E-mail's авторов



\rukeywords{интегральное преобразование Лапласа, интегральные уравнения, обобщённая гипергеометрическая функция}

\ruabstract{Рассматриваются обобщённые интегральные преобразования Лапласа, которые в ядре содержат обобщённую конфлюэнтную гипергеометрическую функцию ${\mathstrut}_1\Phi{\mathstrut}_1^{\tau,\beta}(a;c;z)$. 
С использованием свойств этих преобразований для них получен аналог теоремы о свёртке. 
Методом интегральных преобразований решены интегральные уравнения Вольтерра первого рода, содержащие в ядре конфлюэнтную гипергеометрическую функцию. При решении интегральных уравнений использовались формулы обращения введённых интегральных преобразований, полученные автором ранее.}

\entitle{Generalized Integral Laplace Transform  and Its Application to Solving Some Integral Equations}

\enauthor{S.~M.~Zaikina\inst{1,2}}{S.~M.~Z\,a\,i\,k\,i\,n\,a}


\enaffil{\ensamgtu. \and \envolgogradsu.}

\ruauthordetails{1}{\fullauthor{\rupid{65335}{Светлана Михайловна Заикина}}\position{аспирант}
\fdept{каф. прикладной математики и информатики\inst{1}}\lposition{ассистент}\dept{каф. компьютерных наук и экспериментальной математики\inst{2}}
}

\enauthordetails{1}{\fullauthor{\enpid{65335}{Svetlana M. Zaikina}} 
\position{Postgraduate Student} \fdept{Dept. of Applied Mathematics \&~Computer Science\inst{1}}
\lposition{Assistant}\dept{Dept. of Computer Science \&  and Experimental Mathema\-tics\inst{2}}}


\enabstract{We present integral transforms $\widetilde {\mathcal L}\left\{f(t);x\right\}$ and $\widetilde {\mathcal L}_{\gamma_1,\gamma_2,\gamma} \left\{f(t);x\right\}$, generalizing the classical Laplace transform. The $(\tau, \beta)$- generalized confluent hypergeometric functions are the kernels of these integral transforms. At certain values of the parameters these transforms coincides with the famous classical Laplace transform. The inverse formula for the transforms is given. The convolution theorem for transform $\widetilde {\mathcal L}\left\{f(t);x\right\}$ is proven. Volterra integral equations of the first kind with core containing the generalized confluent hypergeometric function ${\mathstrut}_1\Phi{\mathstrut}_1^{\tau,\beta}(a;c;z)$ are considered. The above equation is solved by the method of integral transforms. The treatment of integral transforms is applied to get the desired solution of the integral equation. The solution is obtained in explicit form.}



\enkeywords{Laplace integral transform, integral equations, generalized hypergeometric function}

\date{30/IX/2013} % Дата поступления статьи в редакцию.
\revisiondate{05/XII/2013} % В окончательном виде
\accepteddate{17/I/2014}

%%%%%%%%%%%%%%%%%%%%%%%%%%%%%%%%%%%%%%%%%%%%%%%%%%%%%%%%%%%%%%%%%%%%%%%%%%%%%%%%%%%%


\makerutitle


Метод интегральных преобразований широко используется для решения дифференциальных и интегральных уравнений, а сама теория интегральных преобразований является одной из важных ветвей прикладного анализа. 
Многие задачи математической физики, астрофизики, термодинамики, механики и других естественных наук приводят к необходимости применения теории интегральных преобразований.
   
Преимущества метода интегральных преобразований заключаются в том, что он даёт возможность: 
\begin{enumerate}
	\item  сведения сложных задач к менее сложным;
	\item  получения окончательного результата в явном виде.
\end{enumerate}
   
Разработкой теории и методов интегральных преобразований занимались Г.~Бейтмен, В.~А.~Диткин,  A.~A.~Килбас, А.~П.~Прудников, С.~М.~Ситник, А.~Эрдейи, M.~Saigo, I.~N.~Sneddon [\citen{zaika:bib:1}--\citen{zaika:bib:4}] и многие другие.
   
Изучение интегральных уравнений первого рода и так называемых парных, тройных и $N$-арных интегральных уравнений, которые часто встречаются в приложениях, приводит к необходимости рассматривать интегральные преобразования со специальными функциями в ядрах. 
Эти вопросы изучены в работах [\citen{zaika:bib:3}, \citen{zaika:bib:5}].
   
   
Рассмотрим обобщённые интегральные преобразования Лапласа:   
\begin{gather}
\label{zaik:1}
\widetilde {\mathcal L}\left\{f(t);x\right\}=\int_0^{\infty}
\exp(-tx) {\mathstrut}_1\Phi {\mathstrut}_1^{\tau,\beta}\left(a; c; -b(tx)^{\gamma}\right)f(t)dt,
\\
\label{zaik:2}
\widetilde {\mathcal L}_{\gamma_1,\gamma_2,\gamma} \left\{f(t);x\right\}=\int_0^{\infty}t^{\gamma_2}
\exp(-(xt)^{\gamma_1}) {\mathstrut}_1\Phi {\mathstrut}_1^{\tau,\beta}\left(a; c; -b(tx)^{\gamma \gamma_1}\right) f(t)dt,
\end{gather}
где $t\geq 0$; 
$\gamma\in \mathbb C$, 
$\gamma_1>0$, $\gamma_2>0$, $b\geq 0$; 
$f(t)\equiv 0$ при $t<0$;  $t^{\gamma_2}f(t)\hm <M \exp(s_0 t^{\gamma_1})$; 
$M$, $s_0$ "--- постоянные числа (при $t>0$);
${\mathstrut}_1\Phi{\mathstrut}_1^{\tau,\beta}$ "--- обобщённая конфлюэнтная гипергеометрическая функция [\citen{zaika:bib:6}]:
\begin{gather*}
{\mathstrut}_1\Phi{\mathstrut}_1^{\tau,\beta}(a,c;z)=
\frac{\Gamma(c)}{\Gamma(a)\Gamma(c-a)} 
\int_0^1 t^{a-1}(1-t)^{c-a-1} {\mathstrut}_1\Psi{\mathstrut}_1 
     \left[ \left. \begin{array}{c}
             (c,\tau)\\
             (c,\beta)
            \end{array} \right| zt^\tau
     \right] dt, 
\\
\mathop{\rm Re} c>\mathop{\rm Re} a>0, \quad  \tau>0, \quad  \tau\in \mathbb R, \quad  \beta>0, \quad  \beta\in \mathbb R, \quad  \tau-\beta<1;\\
\Gamma(z)=\int_0^{\infty}e^{-t}t^{z-1}dt
\end{gather*}
 "---  классическая Гамма-функция;
$${\mathstrut}_1\Psi{\mathstrut}_1 
     \left[ \left. \begin{array}{c}
             (c,\tau)\\
             (c,\beta)
            \end{array} \right| zt^\tau
     \right]=\sum\limits_{n=0}^{\infty}\frac{\Gamma(c+\tau n)}{\Gamma(c+\beta n)}\frac{z^n t^{\tau n}}{n!} $$ 
"--- частный случай функции Райта [\citen{zaika:bib:7}].
     
Заметим, что при $b=0$ преобразование \eqref{zaik:1} совпадает с классическим преобразованием Лапласа [\citen{zaika:bib:1}]:
\begin{equation}
\label{zaik:3}
 \mathcal L\left\{f(t);x\right\}=  \int_0^{\infty}\exp(-tx)f(t)dt.
\end{equation}
Аналогично при $\gamma_2=0$, $\gamma_1=1$, $b=0$ преобразование \eqref{zaik:2} совпадает с \eqref{zaik:3}.
   
В работах  [\citen{zaika:bib:8}, \citen{zaika:bib:9}] приведены и доказаны формулы обращения для интегральных преобразований \eqref{zaik:1}, \eqref{zaik:2}.
Например, для интегрального преобразования \eqref{zaik:1} справедлива формула
\begin{equation}
\label{zaik:4}
f(u)=\frac{\Gamma(a)}{\Gamma(c)}\int_0^{\infty}(ux)^{-1}g(x)\mathcal K(ux)dx,
 \end{equation}
где
\[
\mathcal K(x)=\frac{1}{2\pi i}\int_{\sigma-i\infty}^{\sigma+i\infty}\frac{x^s}{\xi(s)}ds,\quad 
g(y)=\widetilde{\mathcal L}\left\{f(x);y\right\},
\]\[
\xi(s)={\mathstrut}_2\Psi{\mathstrut}_1 \left[ \left. \begin{array}{c}
                    (a,\tau);(s,\gamma)\\
                    (c,\beta)
                    \end{array} \right| -b \right].
\]
 
 %Рассмотрим интегральный оператор, содержащий обобщённую конфлюэнтную гипергеометрическую функцию ${_1}\Phi_1^{\tau,\beta}$
 % \begin{equation}\label{n5}
%\int_0^{x}g(t)f(x-t){_1}\Phi_1^{\tau,\beta}\left(a; c; -b(st)^{\gamma}\right)dt=F(x), 
% \end{equation}
% где $0<x<\infty; Re c>Re a>0, \tau\in \R, \tau>0, \beta\in \R, 
% \beta>0, \tau-\beta<1 $.
 
% Непосредственными вычислениями убеждаемся в справедливости равенства
 
% \begin{equation}\label{n6}
%\mathcal L \left\{F(x);s\right\}=\widetilde{\mathcal L}\left\{g(t);s\right\}\cdot \mathcal L \left\{f(z);s\right\}
% \end{equation}
 
 
%\textbf{Теорема.}

\smallskip

\begin{theorem}[1]При условиях $\mathop{\rm Re} c>\mathop{\rm Re} a>0,$ 
$\tau\in \mathbb R,$ $\tau>0,$ $\beta\in \mathbb R,$ $\beta>0,$ $\tau-\beta<1,$ 
$h(x)\in L(0, +\infty)$
интегральное уравнение Вольтерра первого рода
\begin{equation}
\label{zaik:5}
\int_0^x g(t)(x-t)^{\gamma-1}{\mathstrut}_1\Phi{\mathstrut}_1^{\tau,\beta}\left(a; c; -b(st)^{\gamma}\right)dt=h(x)
\end{equation}
имеет единственное решение
\begin{equation}
\label{zaik:6}
g(t)=\frac{\Gamma(a)}{\Gamma(c)\Gamma(\gamma)}\int_0^{\infty} x^{\gamma-1}t^{-1}\mathcal K(xt)\cdot\mathcal L\left\{h(y);x\right\}dx.
 \end{equation}
\end{theorem}

\smallskip

\begin{proof}	
Применим к обеим частям уравнения \eqref{zaik:5} преобразование Лапласа:
$$
\int_0^{\infty}e^{-xs}\biggl(\int_0^x g(t)(x-t)^{\gamma-1} {\mathstrut}_1\Phi{\mathstrut}_1^{\tau,\beta}\left(a; c; -b(st)^{\gamma}\right)dt\biggr)dx=\mathcal L\left\{h(x);s\right\}. 
$$
%Учитывая равенство (6), получим
Непосредственными вычислениями убеждаемся в справедливости равенства
$$
\mathcal L\left\{h(x);s\right\}=\widetilde{\mathcal L}\left\{g(x);s\right\}\cdot \mathcal L\left\{x^{\gamma-1};s\right\}.
$$	
Учитывая, что [\citen{zaika:bib:1}] 
$$
\mathcal L\left\{x^{\gamma-1};s\right\}=\frac{\Gamma(\gamma)}{s^{\gamma}},
$$ будем иметь
$$
\widetilde{\mathcal L}\left\{g(t);s\right\}=\frac{s^{\gamma}}{\Gamma(\gamma)}\cdot \mathcal L\left\{h(t);s\right\}.
$$

Применяя формулу обращения \eqref{zaik:4} для интегрального преобразования \linebreak $\widetilde{\mathcal L}\left\{g(t);s\right\}$, получим искомую функцию \eqref{zaik:6}.
\end{proof}

\smallskip

\begin{theorem}[2]При условиях существования интегралов
\begin{gather*}
\varphi(s)=\mathcal L_m\left\{f(t);s\right\}=\int_0^{\infty}t^{m-1}e^{-(ts)^m}f(t)dt, 
\\
\psi(s)=\mathcal L_m\left\{g(t);s\right\}=\int_0^{\infty}t^{m-1}e^{-(ts)^m}g(t)dt
\end{gather*}
и их абсолютной сходимости имеет место следующее равенство\/{\rm:}
\begin{equation}
\label{zaik:7} 
   \mathcal L_m\left\{F(x);s\right\}=\mathcal L_m\left\{f(t);s\right\}\cdot \mathcal L_m\left\{g(t);s\right\},
    \end{equation} 
где
    $$F(x)=\int_0^{x}y^{m-1}g(y)f\left(\sqrt[m]{x^m-y^m}\right)dy.$$
\end{theorem}

\smallskip

\begin{mproof}
Здесь $f(t)=0$ и $g(t)=0$ при $t<0$.
При доказательстве формулы \eqref{zaik:7} воспользуемся формулой обращения для интегрального преобразования $\mathcal L_m$:
$$
f(t)=\frac{m}{2\pi i}\int_{\sigma-i\infty}^{\sigma+i\infty}\varphi(s^{{1}/{m}})e^{st^m}ds.
$$

Рассмотрим интеграл
\begin{multline*}
\frac{m}{2\pi i}\int_{\sigma-i\infty}^{\sigma+i\infty}
\varphi(s^{{1}/{m}})\psi(s^{{1}/{m}})e^{st^m}ds=
\\
=\frac{m}{2\pi i}\int_{\sigma-i\infty}^{\sigma+i\infty}
\varphi(s^{{1}/{m}})e^{st^m} \biggl(
\int_{0}^{\infty}y^{m-1}e^{-sy^m}g(y)dy\biggr)ds=
\\
=\int_{0}^{\infty}g(y)y^{m-1}dy
\biggl(
\frac{m}{2\pi i}\int_{\sigma-i\infty}^{\sigma+i\infty}e^{s(t^m-y^m)}\varphi
(s^{{1}/{m}})ds\biggr)=
\\
=\int_{0}^{\infty}g(y)y^{m-1}f\left(\sqrt[m]{t^m-y^m}\right)dy= \int_{0}^{t}y^{m-1}g(y)f\left(\sqrt[m]{t^m-y^m}\right)dy,
\end{multline*}
так как $f(t)=0$ при $t<0$.



Введём обозначение
$$f\ast g=\int_0^t y^{m-1}g(y)f\left(\sqrt[m]{t^m-y^m}\right)dy.$$
Тогда формула \eqref{zaik:7} запишется в виде
$$
\mathcal L_m\left\{f\ast g;s\right\}=\varphi(s)\cdot\psi(s). \quad \square
$$

\end{mproof}

\smallskip

Рассмотрим применение формулы \eqref{zaik:7} для решения интегрального уравнения, содержащего в ядре обобщённую конфлюэнтную гипергеометрическую функцию:
\begin{equation}
\label{zaik:8} 
\int_0^{x} y^{m-1}(x^m-y^m)^{{\gamma}/{m}}{\mathstrut}_1\Phi{\mathstrut}_1^{\tau,\beta}
\bigl(a;c;-b(x^m-y^m)^{{\gamma_1}/{m}}\bigr)\cdot g(y)dy=F(x).
 \end{equation} 
Применим к обеим частям уравнения \eqref{zaik:8} интегральное преобразование $\mathcal L_m$.
 
Поскольку
$$
f\left((x^m-y^m)^{{1}/{m}}\right)=
(x^m-y^m)^{{1}/{m}}\cdot {\mathstrut}_1\Phi{\mathstrut}_1^{\tau,\beta}
\bigl(a;c;-b(x^m-y^m)^{{1}/{m}}\bigr),
$$
в силу равенства \eqref{zaik:7}  получим
\begin{equation}
\label{zaik:9} 
\mathcal L_m\left\{f(z);s\right\}\cdot\mathcal L_m\left\{g(t);s\right\}=\mathcal L_m\left\{F(x);s\right\} ,
 \end{equation}
где $z=(x^m-y^m)^{{1}/{m}}$, то есть  $f(z)=z^{\gamma} {\mathstrut}_1\Phi{\mathstrut}_1^{\tau,\beta}\left(a;c;-bz^{\gamma_1}\right).$

Вычислим
\begin{multline*}
\mathcal L_m\bigl\{z^{\gamma} {_1}\Phi_1^{\tau,\beta}\left(a;c;-bz^{\gamma_1}\right);s\bigr\}=
\int_0^{\infty}z^{m-1}e^{-z^ms^m}z^{\gamma}{\mathstrut}_1\Phi {\mathstrut}_1^{\tau,\beta}\left(a;c;-bz^{\gamma_1}\right)dz=
\\
=
\frac{\Gamma(c)}{\Gamma(a)}\int _0^{\infty}z^{m-1+\gamma}e^{-z^ms^m}\cdot \sum\limits_{n=0}^{\infty}
  \frac{\Gamma(a+\tau n)}{\Gamma(c+\beta n)}\frac{(-bz^{\gamma_1})^n}{n!}dz=
  \\
=
\frac{\Gamma(c)}{\Gamma(a)}\sum_{n=0}^{\infty}
  \frac{\Gamma(a+\tau n)}{\Gamma(c+\beta n)}\frac{(-b)^n}{n!} \int_0^{\infty}e^{-z^ms^m}z^{m-1+\gamma+\gamma_1 n}dz= 
  \\
=\frac{\Gamma(c)}{\Gamma(a)}\frac{1}{ms^{m+\gamma}}\sum\limits_{n=0}^{\infty}
  \frac{\Gamma(a+\tau n)\Gamma\left(\frac{\gamma+\gamma_1 n}{m}+1\right)}{\Gamma(c+\beta n)}\cdot\frac{(-b)^n}{n!s^{\gamma_1 n}}=
\\
=\frac{1}{ms^{m+\gamma}}\frac{\Gamma(c)}{\Gamma(a)}\cdot 
  {\mathstrut}_2\Psi {\mathstrut}_1\left[\left.
  \begin{array}{c}
    (a;\tau); \left(\frac{\gamma}{m}+1;\frac{\gamma_1}{m}\right)\\
    (c;\beta)
  \end{array}\right| \frac{-b}{s^{\gamma_1}}\right].
\end{multline*}
Здесь использовалась замена переменной $ z^ms^m=t$.

Следовательно, уравнение \eqref{zaik:9} можно записать в виде
  $$\frac{1}{ms^{m+\gamma}}\frac{\Gamma(c)}{\Gamma(a)} 
  {\mathstrut}_2\Psi{\mathstrut}_1\left[\left.
  \begin{array}{c}
    (a;\tau); \left(\frac{\gamma}{m}+1;\frac{\gamma_1}{m}\right)\\
    (c;\beta)
  \end{array}\right| \frac{-b}{s^{\gamma_1}}\right]\cdot \mathcal L_m\left\{g(t);s\right\}=
\mathcal L_m\left\{F(x);s\right\}.$$
  Отсюда 
  $$\mathcal L_m\left\{g(t);s\right\}=ms^{m+\gamma}\frac{\Gamma(a)}{\Gamma(c)}\left(
  {\mathstrut}_2\Psi{\mathstrut}_1\left[\left.
  \begin{array}{c}
    (a;\tau); \left(\frac{\gamma}{m}+1;\frac{\gamma_1}{m}\right)\\
    (c;\beta)
  \end{array}\right| \frac{-b}{s^{\gamma_1}}\right]\right)^{-1}\cdot \mathcal L_m\left\{F(x);s\right\}.$$
  Применяя формулу обращения для преобразования $\mathcal L_m$, получим решение уравнения \eqref{zaik:8}:
\begin{multline*}
  g(t)=\frac{m^2}{2\pi i}\frac{\Gamma(a)}{\Gamma(c)}\int_{\sigma-i\infty}^{\sigma+i\infty}
  e^{st^m}s^{({m+\gamma})/{m}}
  \left(  {\mathstrut}_2\Psi{\mathstrut}_1\left[\left.
  \begin{array}{c}
    (a;\tau); \left(\frac{\gamma}{m}+1;\frac{\gamma_1}{m}\right)\\
    (c;\beta)
  \end{array}\right| \frac{-b}{s^{\frac{\gamma_1}{m}}}\right]\right)^{-1}\times \\
\times \mathcal L_m\left\{F(x);s^{{1}/{m}}\right\}ds.
\end{multline*}




\begin{thebibliography}{9}

\Bibitem{zaika:bib:1}
\by A.~Erd\'elyi
\book Tables of Integral Transforms \rm (Bateman Manuscript Project) 
\publ McGraw-Hill
\publaddr New York
\yr 1954
\miscnote vol. 1, Moscow, Nauka,  1969; vol. 2, Moscow, Nauka, 1970 [Russian translation]
\zmath{http://zbmath.org/?q=an:0058.34103|an:0055.36401|an:0184.33704|an:0204.50001}


\RBibitem{zaika:bib:2}
\by В.~А.~Диткин, А.~П.~Прудников
\book Интегральные преобразования и операционное исчисление
\publaddr М.
\publ Наука
\yr 1974
\totalpages 542
\transl 
\by V.~A.~Ditkin, A.~P.~Prudnikov
\book Integral transforms and operational calculus
\publ Pergamon Press
\yr 1965 
\totalpages xi+529
\publaddr  	Oxford, New York 
\serial International series of monographs in pure and applied mathematics
\vol 78

\Bibitem{zaika:bib:3}
\by A.~A.~Kilbas, M.~Saigo
\book  H-Transforms: Theory and Applications
\serial Series on Analytic Methods and Special Functions
\publ CRC Press
\publaddr Boca Raton 
\vol 9
\yr 2004
\totalpages  xii+389
\zmath{http://zbmath.org/?q=an:01907576}



\Bibitem{zaika:bib:4}
\by I.~N.~Sneddon
\book The use of integral transforms
\publaddr New York etc.
\publ McGraw-Hill Book Comp. 
\totalpages xii+539 
\yr 1972
\zmath{http://zbmath.org/?q=an:0237.44001}

\RBibitem{zaika:bib:5}
\by Н.~О.~Вiрченко
\book Парнi ($N$-арнi) iнтегральнi рiвняння
\publaddr Ки\"\i{}в
\publ Задруга
\yr 2009
\totalpages 476
\misctransl
\by N.~O.~Virchenko
\book Parni ($N$-arni) integral'ni rivnyannya \rm [Pairs ($N$-ary) integral equations]
\publaddr Kiev
\publ Zadruga
\yr 2009
\totalpages 476
\lang In~Ukrainian

\Bibitem{zaika:bib:6}
\paper On the generalized confluent hypergeometric function and its applications
\by N.~Virchenko
\jour Fract. Calc. Appl. Anal.
\vol 9
\yr 2006
\issue 2
\pages 101--108
\zmath{http://zbmath.org/?q=an:1128.33006}

\Bibitem{zaika:bib:7}
\paper The Asymptotic Expansion of the Generalized Hypergeometric Function
\by E.~M.~Wright
\jour J.~London Math. Soc.
\yr 1935
\vol s1-10
\issue 4
\pages 286--293.
\crossref{10.1112/jlms/s1-10.40.286}

\RBibitem{zaika:bib:8}
\by О.~А.~Репин, С.~М.~Заикина
\paper Некоторые новые обобщенные интегральные преобразования и~их применение в~теории дифференциальных уравнений
\jour Вестн. Сам. гос. техн. ун-та. Сер. Физ.-мат. науки
\yr 2011
\issue 2(23)
\pages 8--16.
\mathnet{http://mi.mathnet.ru/vsgtu913}
\crossref{10.14498/vsgtu913}
\misctransl
\by O.~A.~Repin, S.~M.~Zaikina
\paper Some new generalized integral transformations and their application in differential equations theory
\jour Vestn. Samar. Gos. Tekhn. Univ. Ser. Fiz.-Mat. Nauki
\yr 2011
\issue 2(23)
\pages 8--16.
\lang In Russian

\Bibitem{zaika:bib:9}
\paper On some generalized integral transforms
\by N.~Virchenko, S.~L.~Kalla, S.~Zaikina
\jour Handronic Journal
\yr 2009
\vol 32
\issue 5
\pages 539--548
\zmath{http://zbmath.org/?q=an:1207.44002}


\end{thebibliography}

\renewcommand{\doihacke}{\linebreak}


\makeentitle

\renewcommand{\baselinestretch}{0.9}

\newpage




%\end{document}
